
Se genero ademas la tabla de síndromes que se encuentra en el archivo \href{https://github.com/AbelCorvalan0/fec-lab2/tree/main/model/outputs/syndromes/error_syndromes.txt}{\textbf{error syndromes.txt}}. Los vectores se presentan de la siguiente forma:

\begin{table}[h!]
    \centering
    \begin{tabular}{|c|c|}
        \hline
        1. Vector recibido     &
        2. Sindrome \\
        \hline
    \end{tabular}
    % \caption{Descripción de las columnas de salida del decodificador.}
    \label{tab:decoder_output}
\end{table}

% La totalidad de los errores que se pueden producir dentro de los 24 bits
% de una palabra recibida en el decoder son generados de la siguiente forma:

% El orden en que se encuentran las variables es el siguiente:

% \begin{table}[h!]
%     \centering
%     \begin{tabular}{|c|c|c|c|c|}
%         \hline
%         1. Vector recibido     &
%         2. Mensaje              &
%         3. Máscara de error     &
%         4. Palabra corregida    &
%         5. Palabra incorregible \\
%         \hline
%     \end{tabular}
%     \caption{Descripción de las columnas de salida del decodificador.}
%     \label{tab:decoder_output}
% \end{table}

% Los valores decodificados se encuentran en el archivo \href{https://github.com/AbelCorvalan0/fec-lab2/blob/main/model/outputs/encoder/encoder_top_testing/encoder_output.txt}{codeword with errors.txt}.